\makeatletter
\def\input@path{{./}{../}{../../}{../../../}{../../../../}{preamble/}{../preamble/}{../../preamble/}{../../../preamble/}}
\makeatother
\input{preamble_exams.tex}
\input{graphs/highlights.tex}
% \GTenHighlightsfalse % Uncomment to disable instructional highlights.
\input{graphs/specific_highlights.tex}
\input{graphs/uniform_linear_probability_distributions.tex}
% Density outlines are neutral; instructional roles use the shared toggle.
\newcommand{\PracticeDensity}[9]{\begingroup\colorlet{blue}{black}\CompactDensityCoordinatesGraph{#1}{#2}{#3}{#4}{#5}{#6}{#7}{#8}{#9}\endgroup}

\hypersetup{pdftitle={Grade 10 Global Economics: Term 3, Lesson 1, C2},pdfsubject={Home practice: Uniform Continuous Models}}
\title{Grade 10 Global Economics\\Term 3, Lesson 1: Home Practice\\C2 --- Uniform Continuous Models}
\author{}
\date{}
\begin{document}
\maketitle
\setcounter{tocdepth}{1}
\setcounter{secnumdepth}{2}
\phantomsection\label{toc}
\tableofcontents
\section*{Evaluated criteria}
\begin{description}
\item[\textbf{C2}] Applies the uniform continuous distribution on an interval \([a,b]\) to calculate probabilities as areas of rectangles in economic and financial contexts.
\end{description}
\textbf{How to practise.} Write both endpoints before drawing a rectangle. Mark any part of a requested interval that lies outside the support. Try each task on paper before opening its solution link. Record the model, working, units and an interpretation; use the arrows to return.
All datasets and economic or financial models in this activity are illustrative, rather than observed market evidence.

\clearpage
\section{Introduction}\label{sec:introduction}
\subsection*{Highlight legend}
\phantomsection\label{tab:highlight-legend}
\begin{tabularx}{\linewidth}{@{}>{\raggedright\arraybackslash}p{.20\linewidth}>{\raggedright\arraybackslash}p{.22\linewidth}>{\raggedright\arraybackslash}X@{}}
\toprule
\textbf{Symbol or notation} & \textbf{Concept} & \textbf{Brief meaning}\\
\midrule
$\GTenHLValue{X}\sim U(a,b)$ & \GTenHLValue{Uniform model} & Equal density on one interval. \\
$[a,b]$ & Support & The allowed interval. \\
$\GTenHLSecond{1/(b-a)}$ & \GTenHLSecond{Density} & Rectangle height. \\
$\HLProbability{P(c<X<d)}$ & \HLProbability{Probability} & Relevant rectangle area. \\
\bottomrule
\end{tabularx}
\subsection{Introductory example 1: Waiting for a scheduled transfer}\label{ex:intro-one}
A transfer is equally likely to complete at any time $\GTenHLValue{T}$ from $2$ to $10$ minutes: $T\sim U(2,10)$. The \GTenHLSecond{density} is $\GTenHLSecond{1/8}$ per minute. The probability of completion from four to six minutes is the rectangle area $2(1/8)=\HLProbability{1/4}$.
\subsection{Introductory example 2: A small currency exchange fee}\label{ex:intro-two}
An illustrative exchange fee $\GTenHLValue{F}\sim U(1,3)$ dollars has density $1/2$ per dollar. Equal dollar-width bands have equal probabilities, although an exact unrounded fee has probability zero. Widening the support while keeping total area one lowers the rectangle height.
\section{Notation Reference and Practice Procedure}\label{sec:notation-and-procedure}
\subsection{Notation reference}\label{tab:notation-reference}
\begin{tabularx}{\linewidth}{@{}>{\raggedright\arraybackslash}p{.16\linewidth}>{\raggedright\arraybackslash}p{.18\linewidth}>{\raggedright\arraybackslash}p{.27\linewidth}>{\raggedright\arraybackslash}X@{}}
\toprule
\textbf{Symbol or notation} & \textbf{Concept} & \textbf{Definition} & \textbf{Meaning in the introductory examples}\\
\midrule
$\GTenHLValue{X}\sim U(a,b)$ & \GTenHLValue{Uniform model} & Continuous distribution with constant density on $[a,b]$ and zero elsewhere. & \textbullet\ \hyperref[ex:intro-one]{Example 1}: $T$ is uniform from 2 to 10 minutes. \\
$[a,b]$ & Support & Lower and upper bounds with $a<b$; width is $b-a$. & \textbullet\ \hyperref[ex:intro-two]{Example 2}: The fee endpoints are 1 and 3 dollars. \\
$\GTenHLSecond{1/(b-a)}$ & \GTenHLSecond{Density} & Height chosen so that $(b-a)/(b-a)=1$; reciprocal measurement units. & \textbullet\ \hyperref[ex:intro-one]{Example 1}: The height is $1/8$ per minute. \\
$\HLProbability{P(c<X<d)}$ & \HLProbability{Probability} & For $ale c<dle b$, $(d-c)/(b-a)$. & \textbullet\ \hyperref[ex:intro-one]{Example 1}: A two-minute band has area $1/4$. \\
\bottomrule
\end{tabularx}
\subsection{Practice procedure}
\begin{enumerate}
\item Select a uniform model only when constant density across the stated interval is assumed.
\item Write the support width $b-a$ and height $1/(b-a)$. Intersect the requested event with the support.
\item Find rectangle area: relevant width times height. For an upper tail use the width from the threshold to $b$.
\item Interpret the probability and compare changed models by both support width and event width; exact endpoints have zero probability.
\end{enumerate}
\section{General Worked Problems}\label{sec:general-worked}
\SectionProblemsTableOfContents{
\SectionProblemLink{prob:sec:general-worked0}{A transfer completion window}\\
\SectionProblemLink{prob:sec:general-worked1}{Changing the range and clipping an event}\\
\SectionProblemLink{prob:sec:general-worked2}{Finding an unknown upper bound}}

\subsection{A transfer completion window}\label{prob:sec:general-worked0}
A bank models its transfer completion time as $T\sim U(2,10)$ minutes.
\QuestionsTasks[sec:general-worked]{{2}/{sol:sec:general-worked0}/{Find the density, the probability of completion between three and seven minutes, and the probability of waiting more than eight minutes. Explain the units.}}
\subsubsection*{Worked solution}\label{sol:sec:general-worked0}
The support width is $10-2=8$ minutes, so $f(t)=\GTenHLSecond{1/8}$ per minute inside the support and zero outside. The requested rectangles have widths four and two:
\[
P(3<T<7)=4(1/8)=\HLProbability{0.50},\qquad
P(T>8)=(10-8)(1/8)=\HLProbability{0.25}.
\]
Thus $50\%$ of transfers take between three and seven minutes and $25\%$ exceed eight minutes. Minute widths cancel the reciprocal-minute density units.
\PracticeDensity{1}{11}{0.18}{2,3,7,8,10}{2,3,7,8,10}{0,0.125}{t}{0,$1/8$}{(2,0) (2,0.125) (10,0.125) (10,0)}
\secInicio[prob:sec:general-worked0][sec:general-worked]


\subsection{Changing the range and clipping an event}\label{prob:sec:general-worked1}
Provider A has a handling fee $F_A\sim U(0,10)$ dollars. Provider B has $F_B\sim U(2,8)$ dollars.
\QuestionsTasks[sec:general-worked]{{2}/{sol:sec:general-worked1}/{Compare the probabilities of a fee from two to seven dollars and of a fee from minus one to three dollars. Explain why a negative fee is not part of either model.}}
\subsubsection*{Worked solution}\label{sol:sec:general-worked1}
Each provider has constant density, but the heights differ: $1/10$ and $1/6$ per dollar.
\[
P(2<F_A<7)=5/10=\HLProbability{0.50},\qquad
P(2<F_B<7)=5/6.
\]
The second event must be restricted to each support. For A its relevant interval is $[0,3]$; for B it is $[2,3]$:
\[
P(-1<F_A<3)=3/10,\qquad P(-1<F_B<3)=1/6.
\]
B is more likely to fall in the first band, whereas A is more likely to fall in the second. Outside-support widths contribute zero area, so a single provider need not be more likely for every event.
\secInicio[prob:sec:general-worked1][sec:general-worked]


\subsection{Finding an unknown upper bound}\label{prob:sec:general-worked2}
A warehouse models dispatch time as $T\sim U(2,b)$ minutes with $b>8$. It states that $P(T\le8)=0.75$.
\QuestionsTasks[sec:general-worked]{{2}/{sol:sec:general-worked2}/{Determine the upper endpoint, then calculate the probability that dispatch takes between nine and ten minutes.}}
\subsubsection*{Worked solution}\label{sol:sec:general-worked2}
The area up to eight uses width $8-2=6$, and the total support has width $b-2$. Thus
\[
\frac{6}{b-2}=0.75,\quad 6=0.75(b-2),\quad b-2=8,\quad b=10.
\]
The height is $1/8$ per minute, giving
\[
P(9<T<10)=1(1/8)=\HLProbability{0.125}.
\]
The assumed upper bound is ten minutes and $12.5\%$ of dispatches lie in the final minute. This conclusion depends on the uniform assumption; the stated probability alone would not determine a general density.
\secInicio[prob:sec:general-worked2][sec:general-worked]
\section{Advanced Problems}\label{sec:advanced}
\SectionProblemsTableOfContents{
\SectionProblemLink{prob:sec:advanced0}{Choosing a service for a deadline}}

\subsection{Choosing a service for a deadline}\label{prob:sec:advanced0}
For a payment service, $T_A\sim U(0,12)$ minutes and $T_B\sim U(3,9)$ minutes. A customer first requires completion by four minutes, then changes the deadline to eight minutes.
\QuestionsTasks[sec:advanced]{{2}/{sol:sec:advanced0}/{Calculate both deadline probabilities for both services and justify how the preferred service changes. Find a deadline strictly between three and nine minutes at which the probabilities are equal.}}
\subsubsection*{Worked solution}\label{sol:sec:advanced0}
At four minutes the rectangle areas are
\[
P(T_A\le4)=4/12=1/3,\qquad P(T_B\le4)=(4-3)/6=1/6.
\]
A has the greater chance of meeting this tight deadline. At eight minutes,
\[
P(T_A\le8)=8/12=2/3,\qquad P(T_B\le8)=(8-3)/6=5/6,
\]
so B is preferable for that deadline. For $3<t<9$, equate the areas:
\[
t/12=(t-3)/6,\quad t=2t-6,\quad t=6\text{ minutes}.
\]
Both give probability $1/2$ at six minutes. B has a narrower support, but its three-minute minimum delay matters for early deadlines; a narrower support alone does not settle the choice.
\secInicio[prob:sec:advanced0][sec:advanced]

\end{document}

